% displayname : Conduction électrique
\documentclass[12pt,fancy]{cours}%$
\Chapit{5}{6}
\begin{document}
\maketitre{
Pourquoi les métaux sont-ils souvent de bons conducteurs électriques, et qu'est ce qui les différencie des matériaux dits isolants comme le bois ? D'où provient la loi d'Ohm si souvent utilisée $U=RI$ ? Ce chapitre s'intéresse au phénomène de transport de charges dans un métal: la conduction électrique. Nous établirons l'expression simplifiée des équations de Maxwell dans un métal et quelques conséquences dans le régime de l'ARQS. Nous terminerons enfin l'étude par un effet intéressant, l'effet Hall, qui permet notamment de mesurer un champ magnétique à partir d'une mesure de différence de potentiel électrique.
}{
  \item
}


%%%%%%%%%%%%%%%%%%%%%%%%%%%%%%%%%%%%%%%%%%%%%%%%%%%%%%%%%%%%%%%%%%%%%%%%%%
\section{Conduction électrique dans les métaux}

\subsection{Modèle de Drude}

\parag{Le conducteur ohmique}
On appelle \imp{conducteur ohmique} un matériau conducteur de l'électricité  composé de charges \imp{libres} mobiles se déplaçant au milieu de particules immobiles. Comme exemple de conducteurs ohmiques, citons les solutions électrolytiques (les charges mobiles sont des ions), les semi-conducteurs (les charges mobiles sont des électrons et des trous), et les métaux (Cu, Ag, Au, ... les charges mobiles sont des électrons). Dans cette section nous étudierons les métaux.
%, ..., charge mobile = électron de conduction).
% \begin{remarque}
%   Certains milieux conducteurs ne sont pas des conducteurs ohmiques, comme le plasma, car toutes les particules du milieu sont chargées et mobiles.
% \end{remarque}

\begin{definition}[Métal]
Un \imp{métal} est un conducteur ohmique constitué :
\begin{liste}
 \item d'\imp{électrons de conduction} (charges libres), de densité volumique de charges $\rho_{\rm{e}}$;
 \item de \imp{cations} formant le réseau cristallin (charges liées), de densité volumique de charge $\rho_{\rm c}$. 
\end{liste}
La neutralité électrique  du métal impose  $\rho_{\rm e}  + \rho_{\rm c} =0$.
 \end{definition}

 \begin{figure}[h!]
 \centering
 \includegraphics[width=0.4\textwidth]{metal_wikipedia.png}
 \caption{Représentation schématique d'un métal avec un réseau cristallin fixe et des électrons de conduction mobiles. Ce sont ces charges mobiles qui sont responsables du courant de conduction $\j=\rho_{\rm{e}}\vec{v}$ avec $\vec{v}$ la vitesse moyenne des électrons. [Wikipédia].}
 \end{figure}

 \begin{remarque}

 Certains conducteurs ne sont pas des \emph{conducteurs ohmiques}. C'est le cas du plasma, où tous les porteurs de charge sont \emph{libres}. Les solutions électrolytiques NaCl(aq) sont des conducteurs ohmiques, mais les particules liées ne sont pas chargé (molécules du solvant). Les charges mobiles sont les ions dissous.\\
Il existe des milieux où tous les porteurs sont \emph{liés} : il ne peut pas exister de courant de conduction. C'est le cas des \emph{isolants électriques} (air, bois, tissu, ...).
 \end{remarque}

\parag{Modèle de Drude}
 \label{subsec:modèle_drude}
Que se passe-t-il à l'échelle microscopique lorsqu'un conducteur ohmique est soumis à un champ électrique ? Une modélisation simple a été proposée par Drude en 1900.

Un métal est modélisé par un ensemble de $n$ électrons de conduction par unité de volume, de masse $m$, de charge électrique $-e$, de densité volumique de charge $\rho_e = - n e$ ; et de cations fixes. Au cours de son mouvement, un électron de conduction, de vitesse $\Vec{v}$, est freiné en raison de son interaction avec les ions du réseau cristallin. On modélise cette action par une \imp{force de friction} phénoménologique proportionnelle à la vitesse.
\begin{prop}[Modèle de Drude]\label{prop:}
L'action du réseau cristallin sur les électrons de conduction est modélisée par une force de friction
  \eqbox{
    \Vec{F} = -\frac{m}{\tau}\vv{v}.
  }
\end{prop}
% \begin{equation}
% \end{equation}
% consiste à supposer qu'en moyenne les électrons de conduction de masse $m$ et de charge $-e$ se déplacent à une vitesse $\vv{v}$ et sont soumis en plus de la force de Lorentz électrique à une force de friction de la part des charges fixes du réseau de la forme $-\frac{m}{\tau}\vv{v}$.
 La grandeur $\tau$ est homogène à un temps et représente la durée typique entre deux collisions successives de l'électron sur un défaut du réseau cristallin. Nous estimerons $\tau$ plus tard dans le cours.
  % Typiquement $\tau\sim\SI{e-14}{s}$.

\subsection{Loi d'Ohm et de Joule locales}
\parag{Loi d'Ohm locale}
 Appliquons le PFD dans le modèle de Drude  à un électron de conduction  soumis à un champ électrique extérieur $\Vec{E}$ supposé \imp{uniforme} et \imp{stationnaire}.


 \underline{Système} : électron de conduction, masse $m$, charge $-e$, vitesse $\Vec{v}$.

 \underline{Référentiel} : référentiel lié au métal,  dans lequel les cations sort fixes.

 \underline{Bilan des forces} : force de friction  $\Vec{F} = -\frac{m}{\tau}\vv{v} $ et force électrostatique $\Vec{F}_e = -e\Vec{E}$.

 \underline{PFD} : on a
 \begin{equation}
 \label{eq:PFD_drude}
 m\frac{\d\vv{v}}{\d t}=-e\vv{E}-\frac{m}{\tau}\vv{v}.
 \end{equation}
On suppose le régime \emph{permanent} atteint : $\frac{\d \vv{v}}{\d t}=0$. L'électron atteint alors une vitesse limite
\begin{equation}
\vv{v}_{\rm{lim}}=-\frac{e\tau}{m}\vv{E}.
\end{equation}
On rappelle que les cations, chargées positivement, sont fixes et donc ne contribuent pas au courant électrique. La densité volumique de courant vaut donc $\j=\rho_e\vv{v} =-ne \vv{v}_{\rm lim}$, c'est-à-dire
%  \begin{equation}
\eqbox{
  \j=\frac{ne^2\tau}{m}\vv{E}.
  }
 %=\sigma\vv{E}
%  \end{equation}
On voit que $\j$ est proportionnel au champ électrique. Ce résultat constitue la loi d'Ohm locale, et se généralise à tout conducteur ohmique.
\begin{theorem}[Loi d'Ohm locale]\label{form:}
Un conducteur ohmique soumis à un champ électrique \imp{uniforme} et \imp{stationnaire} vérifie
$$
\j= \sigma\vv{E}.
$$
où $\sigma$ est appelée \emph{conductivité électrique}, exprimée en $\SI{}{S.m^{-1}}$, avec $\SI{1}{S} \text{ (siemens)} = \SI{1}{\Omega^{-1}}$.
\end{theorem}
% où nous avons défini $\sigma$ comme la conductivité électrique, exprimée en $\SI{}{\Omega^{-1}.m^{-1}}$ ou de façon équivalente\footnote{Il me semble contre-productif de définir des unités comme $\SI{1}{Hz}=\SI{1}{s^{-1}}$ ou ici $\SI{1}{S}=\SI{1}{\ohm^{-1}}$. Mais elles restent utilisées couramment et il est donc bon de les connaître.} en \SI{}{S.m^{-1}} (siemens par mètre).
\noindent
 Dans le cadre du modèle microscopique de Drude, et pour un champ électrique stationnaire (régime continu), la conductivité électrique a l'expression suivante.

 \begin{theorem}[Conductivité électrique]\label{form:}
  La conductivité élecrique d'um métal en régime continu vaut
  $$\sigma = \frac{ne^2\tau}{m}.$$
 \end{theorem}
% \begin{equation}
  % \eqbox{
    % }
% \end{equation}
Quelques ordres de grandeur sont donnés dans le \Tab{conductivite}.
  \begin{table}[H]
    \centering
    \renewcommand{\arraystretch}{1.2}
    \begin{tabular}{|c|c|c|c|c|c|c|}
      \hline
      Conducteur & Ag & Cu & eau de mer & eau pure & air sec & teflon  \\
      \hline
      $\sigma ~(\SI{}{S.m^{-1}})$ &  $\SI{6.3e7}{}$ & $\SI{6.0e7}{}$  & $\SI{4.8}{}$ & $\SI{6e-5}{}$ & $\SI{e-15}{}$ & $\SI{e-25}{}$\\
      \hline
    \end{tabular}
    \caption{Quelques conductivités électriques de conducteurs ohmiques, du meilleur au moins bon conducteur.}
    \label{tab:conductivite}
    
  \end{table}

%  Pour le cuivre, $\sigma\simeq\SI{6e7}{S.m^{-1}}$ alors que pour l'eau pure $\sigma\simeq\SI{6e-5}{S.m^{-1}}$ ou l'air sec $\sigma\simeq\SI{e-15}{S.m^{-1}}$ (isolants électriques).

  \begin{remarque}
  La résistivité électrique est l'inverse de la conductivité électrique, le plus souvent notée $\rho=1/\sigma$ (en $\SI{}{\ohm.m}$). Attention à ne pas confondre avec une densité volumique de charge !
  \end{remarque}

% \begin{remarque}
% Bien veiller à ne pas confondre la résistivité (homogène à des $\SI{}{\ohm.m}$) avec une densité volumique de charges (homogène à des $\SI{}{C.m^{-3}}$). De même ne pas confondre la conductivité électrique $\sigma$ et la densité surfacique de charges, habituellement aussi notée $\sigma$ ! Pour éviter cette confusion potentielle, la conductivité électrique est parfois notée $\gamma$.
% \end{remarque}

\begin{exemple}{Estimation du temps caractéristique de collision}
  Le cuivre est un très bon conducteur électrique, tout en ayant un coût modéré comparé à l'argent. 
\begin{quest}
  \item	La configuration élecronique du cuivre est $\rm{[Cu] = [Ar] 3d}^{10}{\rm 4s}^1$. Sa masse volumique vaut $\mu = \SI{8.96}{g.cm^{-3}}$, et sa masse molaire $M = \SI{63.5}{g.mol^{-1}}$. Estimer la densité volumique d'électrons de conduction $n$ dans le cuivre.
  \item Estimer le temps de collision $\tau$, sachant que la conductivité du cuivre est $\sigma = \SI{6.0e7}{S.m^{-1}}$. On rappelle la masse d'un électron $m=\SI{9.0e-31}{kg}$, avec $e=\SI{1.6e-19}{C}$.
\end{quest}

\end{exemple}
\examplesol{
On a $\tau = \frac{\sigma m}{n e^2} \simeq \SI{e-14}{s}$.
}


%Une inhomogénéité de potentiel électrostatique $V$ engendre une densité volumique de courant électrique $\j$ qui tend à déplacer des charges électriques, de sorte à s'opposer à l'inhomogénéité qui lui a donné naissance (signe moins dans la loi $\j=-\sigma\,\grad{V}$ avec $\sigma>0)$.

\parag{Loi de Joule locale}
On rappelle que la puissance volumique reçue par les porteurs de charge vaut
\begin{equation}
p_{\rm vol} = \j \cdot \Vec{E}.
\end{equation}
Dans un conducteur ohmique, la loi d'Ohm locale donne
\begin{theorem}[Loi de Joule locale]\label{form:}
La puissance volumique $p$ reçue par les porteurs de charge d'un conducteur ohmique vaut
$$
p_{\rm vol} = \sigma E^2 = \frac{j^2}{\sigma}.
$$
\end{theorem}
\noindent Cette grandeur est positive car $\sigma >0$ : la matière reçoit effectivement de la puissance de la part du champ électromagnétique.

%la densité volumique de puissance (en $\SI{}{W.m^{-3}}$) $p_{\rm{J}}=+\j\cdot\vv{E}$. D'après la loi d'Ohm locale, cette puissance s'exprime aussi comme
%\begin{equation}
%p_{\rm{J}}=\sigma||\vv{E}||^2=\frac{||\j||^2}{\sigma}>0.
%\end{equation}
%Il s'agit d'une grandeur toujours positive, puisque $\sigma>0$. Cette énergie reçue par la matière du conducteur est dissipée sous forme de chaleur : c'est l'effet Joule. Cette dissipation peut s'interpréter microscopiquement par les interactions dissipatives entre les électrons de conduction et le réseau cristallin du métal (force de frottement $\vv{f}=-\frac{m}{\tau}\vv{v}$).


\parag{Transport de charge}

La conduction électrique est en fait un phénomène de transport de charges ! En effet, en régime stationnaire, $\vv{E}=-\grad{V}$. Donc dans un conducteur ohmique,
\eqbox{
  \j=-\sigma\,\grad{V}.
  }
On retrouve une loi phénoménologique caractéristique d'un phénomène de transport (comme pour la loi de Fick ou la loi de Fourier). On peut dresser le tableau d'analogie suivant.

\begin{table}[h!]
\renewcommand{\arraystretch}{1.5}
\begin{tabular}{|c|c|c|c|c|}
\hline
Conduction & Inhomogénéité de & Transport de & Loi de conservation & loi phénoménologique  \\
\hline
thermique & température $T$ & énergie $U=CT$ & $\rho c \dr{T}{t} + \dive{\Vec{\jth}} = 0$ & Fourier $\Vec{\jth} = -\lambda \grad T$ \\
\hline
électrique & potentiel $V$ & charge $q$ & $\dr{q}{t} + \dive{\Vec{\jmath}} = 0$ & Ohm $\Vec{\jmath} = -\sigma \grad V$ \\
\hline
\end{tabular}
\caption{Analogie et différence entre phénomènes de transport : conductions thermique et électrique. Contrairement à la conduction thermique, la conduction électrique ne donne pas lieu à de la diffusion parce que la charge n'est pas proportionnelle au potentiel en général.}
\end{table}

\subsection{Lois d'Ohm et de Joule intégales}

\parag{Loi d'Ohm} Supposons que l'on impose une tension $U$ constante aux bornes d'un conducteur ohmique de longueur $L$ et de section $S$. Alors le champ électrique $\vv{E}=E \ex$ est \imp{uniforme} et \imp{stationnaire} le long du conducteur, de longueur $L$ et de section $S$. 
% Nous connaissons la loi d'Ohm sous sa forme plus familière $U=RI$ où $R$ est la résistance électrique du matériau. Nous allons montrer dans ce qui suit qu'il s'agit en réalité d'une forme intégrale (ou globale) de la loi d'Ohm établie précédemment sous forme locale.
% Trouver une forme intégrale de la loi d'Ohm locale.
La tension entre deux points $A$ et $B$ vaut
\begin{equation}
U=V_A-V_B=\int_{A}^{B}-\d V=\int_{A}^{B}-\grad{V}\cdot \Vec{\d \ell} = \int_{A}^{B}\vv{E}\cdot\vv{\d l}
\end{equation}
On en déduit
\begin{equation}
U=EL=\frac{jL}{\sigma}=\frac{IL}{\sigma S}=RI
\end{equation}
puisque $I=\iint_{S}\j\cdot\vv{\d S}=jS$ (la densité de courant électrique est uniforme). On retrouve la loi d'Ohm intégrale vue en électrocinétique. 
% Cette loi se généralise aux régimes quasi-stationnaires (ARQS).
\begin{theorem}[Loi d'Ohm intégrale]
  La tension $U$ aux bornes d'un conducteur ohmique en régime continu vérifie
$$
U = RI.
$$
\end{theorem}
\noindent
La résistance électrique d'un conducteur filiforme de longueur $L$ et de surface $S$ s'écrit
% \begin{equation}
  \eqbox{
    R=\frac{L}{\sigma S}.
  }
% \end{equation}
% On retrouve une expression analogue à celles obtenues dans les phénomènes de transport de particules et d'énergie thermique. C'est logique puisque la conduction électrique est un phénomène de transport de charges étudié ici dans l'ARQS.
\begin{remarque}
Les notions de résistances n'ont de sens qu'en régime stationnaire ou quasi-stationnaire (ARQS).
\end{remarque}


\parag{Loi de Joule} Dans la situation décrite ci-dessus, la puissance totale reçue par les électrons du conducteur vaut
\begin{equation}
\mc{P} = \pa{ \frac{j^2}{\sigma} } \times V = \pa{ \frac{I^2}{S^2 \sigma} } \times LS = \frac{L}{\sigma S} I^2.
\end{equation}
On retrouve la loi de Joule intégrale.
\begin{theorem}[Loi de Joule intégrale]
Lla puissance électrique $\mc{P}$ dissipée dans un conducteur ohmique en régime continu  vérifie
$$
\mc{P} = RI^2.
$$
\end{theorem}


\subsection{Validité des lois d'Ohm et de Joule}


\parag{Régime de validité en fréquence}
On doit pouvoir considérer le champ électrique comme quasi-stationnaire pour que la loi d'Ohm soit valable. Cela revient à négliger le terme inertiel $m\frac{\d\vv{v}}{\d t}$ devant la force de friction $-\frac{m}{\tau}\vv{v}$ dans le PFD. Si $\Vec{E}(M,t)$ oscille la fréquence $f$, on doit avoir
\begin{equation}
\frac{\norme{m\frac{\d\vv{v}}{\d t}}}{\norme{-\frac{m}{\tau}\vv{v}}} \sim f \tau \ll 1 \Rightarrow  f \ll \frac{1}{\tau} \simeq\SI{e14}{Hz}.
\end{equation}
% \begin{equation}
% f\ll\frac{1}{\tau}.
% \end{equation}
pour un bon conducteur comme le cuivre.
\begin{prop}[Validité de la loi d'Ohm]\label{prop:}
  La loi d'Ohm est valable dans l'ARQS, i.e. à des fréquences $f$ négligeables devant $\SI{e14}{Hz}$.
\end{prop}
% \begin{equation}
% $1/\tau $
% \end{equation}
% La loi d'Ohm est donc valable à des fréquences $f$ négligeables devant $\SI{100}{THz}$.
% et on retrouve l'ordre de grandeur évoqué en Sec.~\ref{subsec:modèle_drude}.
% Usuellement nous n'étudions pas le comportement des métaux à des fréquences aussi élevées, et nous pourrons ainsi supposer que la loi d'Ohm s'applique. Autrement dit, même en régime variable, tant que les fréquences n'excèdent pas quelques \SI{}{THz} on pourra considérer que la loi d'Ohm obtenue en régime permanent s'applique.

\begin{exemple}{Conducteur ohmique en régime variable}
Le PFD donne une équation linéaire en la vitesse. On peut donc utiliser la notation complexe pour la résoudre. Notons $\udl{\Vec{E}}(M,t) = \Vec{E_0}(M) e^{i(\omega t -\varphi_M)}$ le champ électrique et $\udl{\Vec{v}}(t) = \Vec{v_0} e^{i\omega t}$ la vitesse. Montrer que l'on peut généraliser la loi d'Ohm en régime variable en introduisant une conductivité complexe $\udl{\sigma}(\omega)$, dont on donnera l'expression.

\end{exemple}
\examplesol{

Le PFD donne
\begin{align*}
 m\frac{\d\udl{\vv{v}}}{\d t}=-e\udl{\vv{E}}-\frac{m}{\tau}\udl{\vv{v}}
 & \Ar mi\omega  \Vec{v_0}  = -e \Vec{E_0}(M) - \frac{m}{\tau} \Vec{v_0} \\
  & \Ar m\pa{i\omega + \frac{1}{\tau}} \Vec{v_0}  = -e \Vec{E_0}(M) \\
&   \Ar \Vec{v_0} = \frac{-e \tau}{m(1+i\omega \tau)} \Vec{E_0}(M).
\end{align*}
Avec $\je = -n e \udl{\Vec{v}}(t) = \Vec{j_0} e^{i\omega t} $, on trouve
\begin{equation*}
\Vec{j_0} = \frac{(e^2n\tau/m)}{1+i\omega \tau}\Vec{E_0}.
\end{equation*}
On peut donc généraliser la loi d'Ohm en régime variable à condition de définir une conductivité complexe
% \begin{equation}
  \eqbox{
    \udl{\sigma} = \frac{\sigma_0}{1+i\omega \tau},
  }
  où $\sigma_0 = \frac{ne^2\tau}{m}$ désigne la conductivité en régime stationnaire.
% \end{equation}
Un conducteur ohmique agit comme un filtre passe-bas.
}

%\begin{remarque}
%Si on ne suppose pas le champ permanent, on peut dans le cas où $\vv{E}(t)$ varie sinusoïdalement comme $\vv{E}=\vv{E_0}\cos(\omega t)$ passer en notation complexe dans l'équation de la dynamique pour en déduire que $\vv{\underline{j}}=\underline{\sigma}\,\vv{\underline{E}}$ avec
%\begin{equation}
%\underline{\sigma}=\frac{\sigma_0}{1+i\omega\tau}=\frac{\sigma_0}{1+(\omega\tau)^2}+i\frac{-\sigma_0 \omega\tau}{1+(\omega\tau)^2}.
%\end{equation}
%Le fait que la conductivité soit complexe traduit le fait qu'il existe un déphasage entre le courant $\j$ et le champ $\vv{E}$. Au premier ordre en $\omega\tau$ (fréquences inférieures au \SI{}{THz}) ce déphasage est nul et alors $\sigma=\sigma(\omega=0)$ noté $\sigma_0$.
%\end{remarque}

\parag{Effet du champ magnétique}
\label{subsec:effet_B}
Le champ magnétique intervient dans la force de Lorentz, mais nous ne l'avons pas inlu dans le PFD. Le rapport entre les composantes magnétique et électrique de la force de Lorentz vaut, en ordre de grandeur,
 \begin{equation}
 \frac{\norme{-e\vv{v}\pv\vv{B}}}{\norme{-e\vv{E}}}\sim\frac{vB}{E}.
 \end{equation}
Nous montrerons plus tard que dans un conducteur ohmique, les champs électrique et magnétique vérifient  $E\sim \frac{cB}{\sqrt{\omega\tau}}$.
%  \begin{remarque}
% Dans le vide où l'équation de propagation est celle de d'Alembert on a $E\sim cB$. Mais comme nous le verrons l'équation de propagation dans un métal n'est pas une équation de d'Alembert mais une équation de diffusion. La relation de dispersion, qui lie $k$ et $\omega$ pour une onde plane progressive monochromatique n'est donc pas linéaire comme le cas de l'équation de d'Alembert. Cela conduit à l'expression mentionnée ci-dessus.
% \end{remarque}
Dans l'ARQS,
 \begin{equation}
   \frac{vB}{E}\sim \frac{v}{c}\sqrt{\omega\tau}.
   \end{equation}
   Or $\sqrt{\omega\tau}\ll 1$ dans l'ARQS, et de plus, $\frac{v}{c}\ll 1$ car les électrons sont non relativistes. On en déduit que la composante magnétique de la force de Lorentz est toujours négligeable dans l'AQRS.
   \begin{exemple}{Vitesse moyenne des électrons de conduction}
    Calculer la vitesse moyenne $v$ d'un électron de conduction dans un métal. On pourra prendre $S=\SI{1}{mm^2}$, $I=\SI{1}{A}$, $n\simeq\SI{e29}{m^{-3}}$ et $e\simeq\SI{e-19}{C}$.
   \end{exemple}
   \examplesol{
    On peut calculer $v$ à partir de $j=I/S=-ne v$.  On trouve $v\simeq\SI{e-4}{m.s^{-1}} \ll c =\SI{3e8}{m.s^{-1}}$.
   }
%
% \subsection{Neutralité d'un métal}
% Nous avons indiqué en introduction qu'un métal se caractérisait par $\rho=0$. Mais que se passe-t-il si on charge le métal avec une densité volumique $\rho_0$ ?
%
% D'après l'équation de conservation de la charge, $\frac{\partial\rho}{\partial t}+\dive{\j}=0$ et de plus d'après la loi d'Ohm locale $\j=\sigma\vv{E}$. En supposant que $\sigma$ est uniforme (indépendant de l'espace), alors $\dive{(\sigma\vv{E})}=\sigma\,\dive{\vv{E}}=\frac{\sigma\rho}{\varepsilon_0}$ en utilisant l'équation de Maxwell-Gauss. Finalement, la densité volumique de charge dans un métal évolue comme $\frac{\partial \rho}{\partial t}+\frac{\sigma}{\varepsilon_0}\rho=0$ qui admet pour solution
% \begin{equation}
%   \rho(t)=\rho_0\exp(-\frac{t}{\tau})\,\overset{t\gg\tau}{\longrightarrow}\, 0
%   \end{equation}
%   avec $\tau=\frac{\varepsilon_0}{\sigma}$ qui vaut par exemple pour le cuivre environ $\frac{\SI{9e-12}{}}{\SI{6e7}{}}\simeq\SI{e-19}{s}$. Au bout de quelques $\tau$, la variation de charge volumique est effacée. Compte tenu de la valeur extrêmement petite de $\tau$, ce phénomène a lieu quasiment instantanément. Le retour à l'électroneutralité est quasiment immédiat: un conducteur reste donc neutre en tout point et il ne peut exister que des densités surfaciques de charges.
%
%   \begin{remarque}
%   C'est ce qui explique pourquoi on modélise les plaques d'un condensateur par une densité surfacique de charges $\sigma$ et non pas une densité volumique $\rho$.
%   \end{remarque}
%
% \subsection{Négliger le courant de déplacement dans un métal}
% Rappelons l'équation de Maxwell-Ampère: $\rot{\vv{B}}=\mu_0(\j+\vv{j_{\rm{D}}})$ avec $\vv{j_{\rm{D}}}=\varepsilon_0\frac{\partial \vv{E}}{\partial t}$. Comparons en ordre de grandeur le courant de conduction et celui de déplacement en utilisant la loi d'Ohm et en se plaçant en régime sinusoïdal
% \begin{equation}
%  \frac{[||\j||]}{[||\vv{j_{\rm{D}}}||]}=\frac{[\sigma ||\vv{E}||]}{[\varepsilon_0 \omega||\vv{E}||]}\sim\frac{\SI{e19}{s^{-1}}}{\omega}.
%  \end{equation}
%  Autrement dit tant qu'on se place à des fréquences $f\ll\SI{e18}{Hz}$ (quasiment toujours), on peut considérer que le courant de conduction domine complètement le courant de déplacement dans un conducteur ohmique. Dans le conducteur, le champ électromagnétique obéit donc aux équations suivantes:
%  \begin{equation}
%  \begin{cases}
% \dive{\vv{E}}=0,\qquad\qquad\rot{\vv{E}}=-\frac{\partial\vv{B}}{\partial t},\\
% \dive{\vv{B}}=0,\qquad\qquad\rot{\vv{B}}=\mu_0\j.
%  \end{cases}
%  \end{equation}
%
%
% \subsection{Expression de la résistance en régime quasi-stationnaire}
% Nous connaissons la loi d'Ohm sous sa forme plus familière $U=RI$ où $R$ est la résistance électrique du matériau. Nous allons montrer dans ce qui suit qu'il s'agit en réalité d'une forme intégrale (ou globale) de la loi d'Ohm établie précédemment sous forme locale.
% \begin{equation}
% \begin{split}
% U=V_A-V_B=\int_{A}^{B}-\d V=\int_{A}^{B}\vv{E}\cdot\vv{\d l}=El=\frac{jl}{\sigma}=\frac{Il}{\sigma S}:=RI
% \end{split}
% \end{equation}
% où nous avons utilisé que $\vv{E}$ est uniforme le long du conducteur, puis que $I=\iint_{S}\j\cdot\vv{\d S}=jS$ en considérant à nouveau une densité volumique de courant uniforme (conséquence de la loi d'Ohm locale $\j=\sigma\vv{E}$ avec $\vv{E}$ uniforme et $\sigma$ également). D'après l'expression précédente, la résistance électrique d'un conducteur filiforme dans l'ARQS s'écrit
% \begin{equation}
% R=\frac{l}{\sigma S}.
% \end{equation}
% On retrouve une expression analogue à celles obtenues dans les phénomènes de transport de particules et d'énergie thermique. C'est logique puisque la conduction électrique est un phénomène de transport de charges étudié ici dans l'ARQS.
%
% \begin{remarque}
% Les notions de résistances thermiques ou hydrauliques vues dans les chapitres précédents n'ont de sens qu'en régime quasi-permanent. Il en est de même ici pour la résistance électrique, c'est pourquoi il est essentiel de se placer dans l'ARQS.
% \end{remarque}

\section{Effet Hall}
Nous avons vu que pour un onde électromagnétique, le champ magnétique était négligeable devant le champ électrique dans un conducteur ohmique. Mais que se passe-t-il en l'absence de champ électrique extérieur, et en présence seulement d'un champ magnétostatique créé par un aimant ? %C'est le cas de l'effet Hall.

\subsection{Approche descriptive}

\parag{Tension de Hall}
Soit un métal parcouru par un courant longitudinal (selon $\vv{e_x}$) d'intensité $I$, et soumis à un champ magnétique vertical (selon $\vv{e_z}$). On observe expérimentalement l'apparition d'une tension transverse $U_H$ (selon $\ey$) proportionnelle au champ magnétique $B$, à l'intensité du courant $I$, et à l'inverse de l'épaisseur $b$, appelée tension de Hall,
\begin{theorem}[Tension de Hall]\label{form:}
  \begin{equation*}
    U_H  = R_H \frac{IB}{b}
  \end{equation*}
  La constante de Hall $R_H$ ne dépend que de la nature du matériau.
\end{theorem}

\begin{figure}[h!]
\centering
\includegraphics[width=\textwidth]{effet_hall.pdf}
\caption{Effet Hall (vue de dessus à droite). Les flèches sont orientées pour des électrons, de charges $q=-e<0$. La composante magnétique de la force de Lorentz induit un déplacement des charges dans la direction transverse: la plaque vers l'avant reçoit les charges (ici les électrons) et se charge donc négativement. Il apparaît donc un champ de Hall, $\vec{E}_{\rm{H}}$ dirigé de la plaque chargée positivement vers celle chargée négativement, autrement dit une différence de potentiel égale à $V_2-V_1=U_{\rm{H}}=E_{\rm{H}}a$.}
\label{fig:effet_hall}
\end{figure}



\parag{Application à une mesure de champ magnétique}

% Quel est l'intérêt de l'effet Hall ? Nous avons vu que $U_{\rm{H}}=E_{\rm{H}}a$. Or à l'équilibre, on peut relier le champ électrostatique de Hall à la norme du champ magnétique transverse $B$. En effet, $0=-e\vv{v}\pv\vv{B}-e\vv{E_{\rm{H}}}$ et en projection selon $\vv{e_y}$ il vient $vB=E_{\rm{H}}=U_{\rm{H}}/a$. Or $j=I/S=-nev$ avec $S=ab$, donc $v=\frac{I}{-abne}$
% et finalement
% \begin{equation}
% U_{\rm{H}}=\frac{-1}{ne}\frac{IB}{b}.
% \end{equation}
Un intérêt de l'effet Hall est de pouvoir mesurer un champ magnétique $B$ à partir d'une tension $U_H$. C'est le principe des capteurs à effet Hall qui équippent la plupart des teslamètres. De plus, la relation entre $U_H$ et $B$ est linéaire.
%Cette relation montre qu'en connaissant l'intensité du courant $I$ qui traverse le métal, ses propriétés géométriques (la longueur $b$) et physiques (la densité de porteurs de charges mobiles $n$), en mesurant la tension de Hall transverse $U_{\rm{H}}$ on en déduit la valeur du champ magnétique $B$. C'est le principe d'un teslamètre: il utilise des sondes de Hall dans différentes directions pour mesurer les composantes de $\vv{B}$ par une simple mesure de tension en imposant le courant.
En pratique, les sondes à effet Hall utilisent non pas des métaux mais des semi-conducteurs, car leur constante de Hall est bien plus grande, ce qui augmente la tension de Hall et ainsi autorise des mesures de champs magnétiques même peu intenses. Pour la même raison, ces sondes utilisent un ruban plat (pour minimiser $b$).

\subsection{Interprétation de l'effet Hall}


On suppose que les porteurs de charges sont des électrons, de charge $-e$ et de densité particulaire~$n$, de vitesse $\Vec{v} = -v \ex$. Ce déplacement de charge crée un courant de densité $\j = -ne \Vec{v} = nev \ex$ et d'intensité
\begin{equation}
  I = \int_S \j \ps \dS = nevab.
\end{equation}

En établissant le champ magnétique $\B = B \ez$, les électrons subissent la composante magnétique de la force de Lorentz
\begin{equation}
  \vv{F}_L=-e\vv{v}\pv\vv{B} = -evB \ey = -\frac{IB}{nab} \ey.
\end{equation}
Les électrons s'accumulent donc sur la plaque avant du conducteur (côté 1), qui se charge alors négativement. En comparaison, la plaque située à l'arrière du conducteur (côté 2) se charge positivement pour maintenir la neutralité du conducteur. Il apparaît ainsi un champ électrique $\vv{E}_{\rm{H}} = E_H \ey$, appelé champ de Hall, qui exerce une force électrique
\begin{equation}
  \vv{F}_{\rm{H}}=-e\vv{E}_{\rm{H}} = eE_H \ey.
\end{equation}
À l'équilibre, cette force électrique de Hall compense la force magnétique de Lorentz, de sorte que
\begin{equation}
    \vv{F}_{\rm{H}} + \vv{F}_L = \Vec{0} \Rar E_H = \frac{IB}{neab}.
\end{equation}
 Il apparaît donc une tension de Hall $U_H$ dans la direction transverse au courant et au champ magnétique donnée par
\begin{equation}
  U_{\rm{H}}=V_2-V_1=\int_{V_1}^{V_2}\d V=-\int_{1}^{2}\vv{E}_{\rm{H}}\cdot\vv{\d l} = E_{\rm{H}}a = \frac{1}{ne} \frac{IB}{b}.
\end{equation}
On prédit ainsi l'expression de la constante de Hall
% \begin{equation}
  \eqbox[]{
    R_H = \frac{1}{ne}.
  }
% \end{equation}
\begin{exemple}{Ordre de grandeur de la tension de Hall}
  Dans les semi-conducteurs, la densité de porteurs de charges $n$ est beaucoup plus faible que dans les métaux, ce qui explique leur plus grande constante de Hall. Dans un métal, $n \sim \SI{e29}{m^{-3}}$, alors que $n \sim \SI{e19}{m^{-3}}$ pour le germanium. Estimer la tension de Hall pour ces deux matériaux pour une épaisseur $b = \SI{1}{mm}$, une intensité $I=\SI{1}{A}$ et une champ magnétique $B=\SI{1}{T}$. 
\end{exemple}
\examplesol{
On trouve $U_H = \SI{0.06}{nV}$ pour un métal et $U_H = \SI{0.6}{V}$ pour du germanium.
}
% \begin{remarquetech}
  \begin{remarque}
    Ce modèle est valable pour certains métaux comme le cuivre ou l'or. En revanche, il conduit à des résultats aberrants (même le signe n'est pas correct!) pour le fer ou l'aluminium par exemple. Seule une théorie quantique permet d'expliquer complètement l'effet Hall.
  \end{remarque}
% \end{remarquetech}
%
%
% % \subsection{Force de Laplace}
% Un conducteur est constitué de charges fixes, de densité volumique de charges $\rho_{\rm{f}}$ et de charges mobiles de densité volumique de charges $\rho_{\rm{m}}$. Or seules les charges mobiles subissent la force magnétique de Lorentz. Les charges fixes ne subissent-elles aucune force ? Si, puisque d'après ce que nous venons de voir, le champ magnétique induit un champ électrique de Hall $\vv{E}_{\rm{H}}=-\vv{v}\pv\vv{B}$ qui lui agit sur toutes les charges, même celles qui sont fixes. Un volume mésoscopique $\d V$ de conducteur subit donc la force élémentaire
% \begin{equation}
% \vv{\d F}=\rho_{\rm{f}}\d V\vv{E}_{\rm{H}}+\rho_{\rm{m}}\d V\left(\vv{E}_{\rm{H}}+\vv{v}\pv\vv{B}\right).
% \end{equation}
% Comme dans son ensemble, le conducteur est neutre électriquement, $\rho_{\rm{f}}=-\rho_{\rm{m}}$. L'expression précédente se simplifie donc en
% \begin{equation}
% \label{eq:force_laplace}
% \vv{\d F}=\rho_{\rm{m}}\d V\vv{v}\pv\vv{B}=\j\d V \pv\vv{B}.
% \end{equation}
% Tout se passe comme si seules les charges mobiles subissaient la composante magnétique de la force de Lorentz. L'expression \eqref{eq:force_laplace} est l'expression générale de la force électromagnétique exercée sur une portion de conducteur de volume $\d V$, appelée force de Laplace. Pour un fil, nous avons $\j\d V=jS\vv{\d l}=I\vv{\d l}$ et donc
% \begin{equation}
% \vv{\d F}=I\vv{\d l}\pv\vv{B},
% \end{equation}
% on retrouve l'expression de la force de Laplace vue en première année.

\end{document}
